% LaTeX counterpart of the ArXivArticle stylesheet sample.
\documentclass[11pt]{article}
\usepackage[T1]{fontenc}
\usepackage{lmodern}
\usepackage{amsmath,amssymb,amsthm}

\newtheorem{theorem}{Theorem}[section]
\newtheorem{definition}[theorem]{Definition}
\theoremstyle{remark}
\newtheorem{remark}[theorem]{Remark}

\title{A Sample Article}
\author{MathNotebook}
\date{Set with the \texttt{article} class, in Latin Modern.}

\begin{document}

\maketitle

\begin{abstract}
The same document as the ArXivArticle notebook sample, typeset in LaTeX, so the
two can be compared side by side: title block, sectioning, theorem environments
numbered per section, displayed equations, and a bibliography.
\end{abstract}

\section{Statement}

\begin{definition}
A graph is a pair $(V,E)$ with $E$ a set of unordered pairs of distinct elements of $V$.
\end{definition}

\begin{theorem}
The theorem environments share one counter, so this theorem follows the definition
above and both reset at the next section.
\end{theorem}

\begin{equation}
  \sum_{v \in V} \deg(v) = 2\lvert E \rvert
\end{equation}

\begin{proof}
Summing degrees counts every edge twice.
\end{proof}

\section{Displayed mathematics}

An unnumbered display, for an aside:
\[
  \chi(G) \le \Delta(G) + 1
\]

Inline mathematics sits in text at the surrounding size, as in $e^{i\pi} = -1$, and a
numbered display is referenced by the live number the palette copies, as in~\eqref{eq:forms}.

A numbered display, aligned as an \texttt{align} environment converts:
\begin{align}
  \partial \omega &= 0 \label{eq:forms} \\
  \omega \wedge \omega &= -\, \omega \wedge \omega
\end{align}

\section{Subsections}

\subsection{A subsection}
Subsection and subsubsection headings are numbered within their parent and reset with it.

\subsubsection{A subsubsection}

\begin{remark}
Remark, Question, and Observation are set like this one --- italic label, upright body.
\end{remark}

\begin{thebibliography}{9}
\bibitem{matex} S.~Horvat, \emph{MaTeX: LaTeX labels in Mathematica}, \texttt{https://github.com/szhorvat/MaTeX}
\bibitem{ollivier} Y.~Ollivier, \emph{Ricci curvature of Markov chains on metric spaces}, \texttt{https://arxiv.org/abs/0810.5625}
\end{thebibliography}

\end{document}
